\subsection{An embedded Tate interface at the missing corner}

\begin{proposition}[B.4.20; Proved; the other Kato axis]
\label{push:B420}
In the proper smooth Kato family $\mathcal K\to\Delta_s$
of B.4.10 there is a second relative effective Cartier
divisor $\mathcal D\subset\mathcal K$, disjoint from
the persistent elliptic divisor $\mathcal E$.
The surface $\mathcal D$ is smooth and proper over
$\Delta_s$. Its general fibre and central fibre are
\[
 \mathcal D_s=\mathbb C^*/(\kappa(s)s^n)^\mathbb Z
 \quad(s\ne0),\qquad
 \mathcal D_0=\hbox{the reduced rational }n\hbox{-cycle}.
\]
At its nodes the relative local equation is $uv=s$.
Its normal line in $\mathcal K$ is the multidegree-zero
Tate line whose factor on the closing collar is $a(s)$.
In particular it has degree zero on every component
of the central cycle. Scaling the first blowdown
coordinate gives a holomorphic $\mathbb C^*$ action
on the pair $(\mathcal K,\mathcal D)$ over the base.

For $n=12$ this constructs an embedded proper copy of
the actual cusp's elliptic root interface, with its
degree-zero normal line, after choosing the units
$\kappa,a$ to be the actual Tate and normal multipliers.
In the additive period convention $\mathbf u=(\tau,\mu)$,
$W=e^{-2\pi i\zeta_2}$, these multipliers are
$\kappa(s)s^{12}=e^{2\pi i\tau}$ and
$a(s)=e^{-2\pi i\mu}$. The elliptic translation by
$6P$ is realized by first-coordinate scaling
$a(s)^{-6}$.

This is an embedded Cartier and normal-line comparison,
not an identification of relative tubular neighbourhoods.
No lift of the whole second period, attachment to the
global partial filling, or origin completion is asserted.
\end{proposition}

\begin{proof}
Use the modification, embedding and compact collar
domains of B.4.10. Write $w$ for the second blowdown
coordinate and set, on the modified fundamental region,
\[
 h_s=\pi_n^*w-s.
\]
The collar identification is $\sigma_s\pi_n$, where
$\sigma_s(z,w)=(\kappa(s)z,a(s)w+(1-a(s))s)$ in
the last chart. Consequently
\[
 h_s\circ(\sigma_s\pi_n)=a(s)h_s.
 \tag{B.4.20.1}
\]
Its zeros therefore glue to a relative effective
Cartier divisor, and (B.4.20.1) is simultaneously the
normal-line transition. The divisor is closed in the
proper family $\mathcal K\to\Delta_s$, so it is proper.

For $s\ne0$, the line $w=s$ misses every blow-up centre.
In the coordinate $v=w-s$ the contracting germ is
\[
 (x,v)\longmapsto
       (\kappa(s)x(a(s)v+s)^n,a(s)v).
\]
Its restriction to $v=0$ is
$x\mapsto\kappa(s)s^n x$. The generic developing
map of B.4.10 restricts to the identity in this
coordinate, since its product $G_s(v)$ has
$G_s(0)=1$. Thus this is precisely the unmodified
other Hopf axis, with the stated Tate multiplier
and normal multiplier $a(s)$.

The divisor of $\pi_n^*w$ is reduced. After the first
blow-up, both the strict horizontal axis and the
exceptional curve have coefficient one. Every later
centre is the intersection of the newest exceptional
curve with the strict vertical axis. It lies on no
other component of the divisor of $\pi_n^*w$.
The new exceptional coefficient is therefore again
one. Before the end collars are identified, the
divisor is a chain comprising the strict horizontal
axis and the $n$ exceptional pieces. The end annuli
join the two endpoint pieces, leaving the familiar
rational $n$-cycle of the parabolic surface.
It is disjoint from its elliptic curve, as also
described in the parabolic model of
\cite[Example 9.10]{IstratiEtAl2022Root}.

At a crossing of two adjacent pieces, ordinary
blow-up coordinates give $\pi_n^*w=uv$ up to a
holomorphic unit, which can be absorbed in $u$.
Hence $\mathcal D$ has equation $uv=s$. It is
smooth, and its central fibre is reduced nodal.
Away from these nodes its defining equation is
a coordinate minus $s$. This proves smoothness
of the whole divisor and semistability over the disc.
The monomial transition charts are those of the
unit-interval Tate model, and the closing on the
nonzero axis is exactly $x\mapsto\kappa(s)s^n x$.
They identify the marked family with that Tate
model, not only its individual smooth fibres.

The same local defining functions trivialize the
normal line. Their only nontrivial closing factor
is (B.4.20.1), namely the holomorphic unit $a(s)$.
On each rational component the line has degree
zero; equivalently its two adjacent-node degree
corrections cancel its self-intersection $-2$.
For $n=1$ the same statement uses the two branches
of the self-node. The factor $a(s)$ gives the
retained gluing parameter around the cycle.
This is exactly the multidegree-zero extension
of the line with factor $a(s)$ on the smooth
Tate curves.

The first-coordinate scaling commutes with the
modification and the collar maps and preserves
$h_s$. One may equivalently descend the vector
field $x\partial_x$ and its local flow. Properness
over the base makes that flow complete on compact
fibres, uniformly on compact parameter sets;
its $2\pi i$ period yields the stated holomorphic
$\mathbb C^*$ action. On the Tate axis it is
ordinary multiplication of $x$, so a holomorphic
unit $a(s)^{-6}$ realizes the marked translation
by $6P$.

For the actual cusp, the first lattice period
acts in multiplicative coordinates by
\[
 (Z,W)\longmapsto
       (e^{2\pi i\tau}Z,e^{-2\pi i\mu}W).
\]
After the twelfth root change, its first multiplier
is $\kappa(s)s^{12}$ for a holomorphic unit.
The second multiplier tends to the contracting
inverse of $-(5+2\sqrt6)$, so after shrinking
$0<|a(s)|<1$ as required in B.4.10.
The preceding exact Tate charts and the normal
factor identify this interface and its line with
the actual ones. The two section normals in the
ruled model are inverse; one chooses the matching
branch accordingly. More precisely, in the total
line-coordinate convention, this $W$ normal is
$M^{-1}$ for $M=\mathcal O(P-O)$: the theta ratio
with divisor $P-O$ has multiplier $e^{2\pi i\mu}$,
the inverse of $a$. The opposite ruled section
uses $M$. The separate first-coordinate factor
$a^{-6}$ is the translation by $6P$.
These marked quantities are not conflated.

We record the neighbourhood limitation explicitly.
At a central node, the smooth Kato total space
has coordinates $(u,v,s)$ with divisor
$h=uv-s=0$. In coordinates $(u,v,h)$ its map to
the base is
\[
 s=uv-h.
 \tag{B.4.20.2}
\]
It is a submersion. The total space of the normal
line, with the inherited map from $\mathcal D$,
would instead have $s=uv$, with a critical normal
line over that node. Therefore there is no
relative linear tubular-neighbourhood identification
across the nodes. This is a literal local
singularity distinction, not only a possible
global obstruction to linearization.

An inverse-normal plumbing variable $f$ would
give the local equation
\[
 r=fh=f(uv-s).
 \tag{B.4.20.3}
\]
The total fourfold in these coordinates is
smooth and its $r$-central divisor is reduced
normal crossing. For $r\ne0$ one has $f\ne0$
and $s=uv-r/f$, whose derivative in $f$ is
$r/f^2\ne0$. Thus this local plumbing changes
the behaviour at the Tate node; (B.4.20.3) by
itself does not identify the prescribed two-period
Mumford cusp. Global plumbing data and that
period comparison are still missing.

There is also a concrete obstruction to the most
direct global use of this plumbing. For fixed
$s\ne0$, all the $n$ blow-up centres lie on the
first Hopf axis $x=0$, at $[v=-s]$ and its
infinitely near strict-axis points. The other
axis $\mathcal D_s$ is $v=0$, $x\ne0$.
The axes are disjoint, and hence every exceptional
rational curve misses $\mathcal D_s$. A plumbing
that leaves the complement of a small tubular
neighbourhood of $\mathcal D_s$ unchanged retains
these compact curves in its $r\ne0$ fibres over
the same $s$. A complex torus has no rational
curve: a holomorphic map from $\mathbb P^1$
lifts to its vector-space universal cover and
is constant. Thus this fixed-base localized
plumbing of the entire Kato surface cannot
produce the prescribed ordinary two-tori.
Further deformations or a different map to the
original base are outside this obstruction.

The Cartier factor, Tate multiplier, reduced
divisor multiplicities, scaling action and the
two local base-map comparisons are checked in
\path{code/ab_push_tate_interface.py} and
its saved output. Properness rests on the
explicit compact domains of B.4.10.
\end{proof}
